\title{Geometric Interpretation of Voltage Constraints in the Extended Current and Excitation Space}
\author{}
\date{24 September 2026}
\newcommand{\PaperAbstract}{%
An electrically excited synchronous machine has three current coordinates but
only two stator-voltage components.  This dimensional imbalance changes the
geometry fundamentally: in the linear steady-state model the stator-voltage
quadratic form is positive semidefinite with rank two, so its boundary is an
oblique elliptic cylinder rather than an ellipsoid.  The result is derived by
rank, eigenvalue, null-space, and singular-value arguments and is visualized as
a family of congruent fixed-excitation ellipses.  Torque level sets are
classified as hyperbolic cylinders, whereas a positive-resistance copper-loss
budget is a genuine ellipsoid.  Minimum-loss operation becomes a generalized
eigenvalue problem.  The voltage-only maximum-torque problem is shown to be
generically unbounded, and compact thermal or current limits are introduced to
obtain a well-posed high-speed optimization.  Numerical level-set methods and
hidden finite-domain constraints are discussed for measured and nonlinear
machine maps.}
